\subsection{Model uncertainty and the economic margin for randomization}\label{sec:robustmargin}
The certificates can be translated into a robustness statement without treating designed numerical primitives as estimated data. Fix a nominal model, initial law, and policy class. Suppose a nonempty set of alternative models has verified uniform bounds
\begin{equation}\label{eq:modelerrors}
 |J_t^{p,M}-J_t^{p,0}|\leq a,\quad |V_t^M-V_t^0|\leq a,
 \quad |C_0^{p,M}-C_0^{p,0}|\leq b
\end{equation}
for every common policy $p$, state, and date, and every $M$ in the set. Superscript 0 denotes the nominal model. The error recursion in Section~\ref{sec:transfer} supplies such bounds from uniform reward, cost, terminal, and kernel enclosures on a common state space; no reset aggregation is needed for this fixed-space perturbation statement. Uniformity over policies is essential. A fitted residual checked only under the proposed policy does not establish \eqref{eq:modelerrors}.

\begin{proposition}[Portable operating and cost intervals]\label{prop:robust}
Let $\varepsilon>2a$. A nominal policy feasible at allowance $\varepsilon-2a$ is feasible at allowance $\varepsilon$ in every alternative model. If $L_+$ is a nominal lower bound at allowance $\varepsilon+2a$ and $U_-$ is the nominal cost of such a tightened feasible policy, then every model-specific common-policy optimum belongs to
\[
 [\max\{0,L_+-b\},\ U_-+b].
\]
The same interval encloses the optimum that minimizes worst-case implementation cost over policies feasible in all the alternative models. In the latter assertion the upper endpoint uses a single policy, not a different policy for each model.
\end{proposition}
\begin{proof}
For the proposed policy, $J^{p,M}\geq J^{p,0}-a\geq V^0-\varepsilon+a\geq V^M-\varepsilon$. Conversely a policy feasible in a model $M$ satisfies $J^{p,0}\geq J^{p,M}-a\geq V^M-\varepsilon-a\geq V^0-\varepsilon-2a$. Its nominal cost is at least $L_+$ and its cost in $M$ at least $L_+-b$. Take infima for the model-specific statement. For the robust problem, the first implication supplies one policy for all models and the second applies to every robustly feasible policy in each model. Taking the worst-case cost preserves the bounds.
\end{proof}

A strict randomized separation also has an economically interpretable administration-cost allowance. Let $U_R$ be the certified cost of a feasible common lottery and $L_D$ a lower bound for every feasible deterministic policy, for the same operating constraint and initial law. Suppose operating payoffs and dynamics are unchanged by administering the lottery. A nonnegative extra administration charge at date $t$ bounded by $h_t$ increases its expected discounted cost by at most $H=\sum_{t<T}\beta^th_t$. If
\begin{equation}\label{eq:overhead}
 U_R+H<L_D,
\end{equation}
then the lottery remains cheaper than every deterministic feasible policy even after this additional charge. This conclusion does not require administration at every visited state; charging the bound everywhere is conservative. Deterministic rules are assumed not to pay the randomization-specific charge.

For a common per-date ceiling $h$, the certified break-even threshold is
\begin{equation}\label{eq:break_even}
 h_*=(L_D-U_R)_+\bigg/\sum_{t=0}^{T-1}\beta^t.
\end{equation}
Every $h<h_*$ preserves the strict ranking. Equality at $h_*$ yields no strict conclusion. For model uncertainty, use a deterministic nominal lower bound at $\varepsilon+2a$ and a nominal lottery feasible at $\varepsilon-2a$; then replace the numerator by $(L_{D,+}-U_{R,-}-2b)_+$. The current development results report $a=b=0$: they calculate economic margins under known designed primitives, not empirical uncertainty radii.

\begin{table}[t]
\caption{Certified allowance for a per-date randomization administration charge}\label{tab:overhead}
\centering\small
\begin{tabular}{rlrrrr}
\toprule $T$ & Installed rule & $\varepsilon$ & Det. lower & Lottery upper & $h_*$ lower bound\\\midrule
\tablerows{revisions/2026-09-25-r40/paper/generated/overhead_selected.tex}
\bottomrule
\end{tabular}
\par\smallskip\footnotesize Uniform initial law, unchanged development economy. Endpoints and the threshold are rounded outward. A positive threshold means that every smaller constant per-date overhead preserves a strict cost advantage over every feasible deterministic policy. All 126 rule/horizon/tolerance/initial-law margins are retained in the machine-readable results, including zeros.
\end{table}
